% ==============================================================================
% Κεντρικό Αρχείο Παρουσίασης (Master Beamer Presentation)
% Τμήμα Μηχανικών Η/Υ & Πληροφορικής (CEID), Πανεπιστήμιο Πατρών
% ==============================================================================

\documentclass[10pt,aspectratio=169,compress]{beamer}

% Φόρτωση Προοιμίου & Μεταδεδομένων
% ==============================================================================
% Ρυθμίσεις Γλώσσας και Γραμματοσειρών (XeLaTeX)
% ==============================================================================
\usepackage{fontspec}
\usepackage{polyglossia}
\setmainlanguage{greek}
\setotherlanguage{english}

% Modern sans-serif font stack (Fira — designed for Metropolis)
\setsansfont{Fira Sans Light}[
  BoldFont = {Fira Sans SemiBold},
  ItalicFont = {Fira Sans Light Italic},
  BoldItalicFont = {Fira Sans SemiBold Italic},
]
\setmonofont{Fira Mono}[Scale=MatchLowercase]
\newfontfamily\greekfonttt{DejaVu Sans Mono}[Scale=MatchLowercase]

% ==============================================================================
% Beamer Theme — Metropolis (modern, minimal, light)
% ==============================================================================
\usetheme[
  progressbar=frametitle,   % thin progress bar under frame titles
  numbering=fraction,       % slide x / N
  sectionpage=progressbar,  % section divider with progress
  block=fill,               % filled blocks
]{metropolis}

% Χρωματική παλέτα — light & airy CEID tones
\definecolor{ceidblue}{RGB}{0, 71, 133}      % Refined UPatras Blue
\definecolor{accentblue}{RGB}{24, 116, 205}  % Lighter accent
\definecolor{softgray}{RGB}{248, 249, 251}   % Near-white background tint
\definecolor{darktext}{RGB}{35, 35, 45}      % Soft black for text
\definecolor{mGreen}{RGB}{14, 159, 110}      % Metropolis green
\definecolor{mYellow}{RGB}{237, 177, 32}     % Metropolis yellow

% Apply palette
\setbeamercolor{normal text}{fg=darktext, bg=white}
\setbeamercolor{alerted text}{fg=ceidblue}
\setbeamercolor{frametitle}{bg=ceidblue, fg=white}
\setbeamercolor{title separator}{fg=accentblue}
\setbeamercolor{progress bar}{fg=accentblue, bg=accentblue!20}
\setbeamercolor{block title}{fg=white, bg=ceidblue}
\setbeamercolor{block body}{fg=darktext, bg=softgray}
\setbeamercolor{block title alerted}{fg=white, bg=accentblue!85}
\setbeamercolor{block body alerted}{fg=darktext, bg=softgray}
\setbeamercolor{block title example}{fg=white, bg=mGreen!85}
\setbeamercolor{block body example}{fg=darktext, bg=softgray}

% Lightweight item markers & tighter spacing
\setbeamertemplate{itemize items}[circle]

% ==============================================================================
% Πρόσθετα Πακέτα
% ==============================================================================
\usepackage{amsmath,amssymb,bm}
\usepackage{graphicx}
\graphicspath{{figures/}{../figures/}{../../evaluation/figures/}}
\usepackage{booktabs}
\usepackage{tabularx}
\usepackage{tikz}
\usetikzlibrary{shapes,arrows,positioning,calc}
\usepackage{appendixnumberbeamer}

% ==============================================================================
% Μεταδεδομένα Τίτλου Παρουσίασης
% ==============================================================================
\title[Distributed GA \& LEAD DHT for UAVs]{Κατανεμημένος Γενετικός Αλγόριθμος για τη Βελτιστοποίηση Αεροδυναμικής UAV}
\subtitle{Χρήση LEAD DHT και Multi-Tier Surrogate (Υποστήριξη Διπλωματικής)}
\author[Ε. Π. Χριστοδουλόπουλος]{Ευστάθιος Παναγιώτης Χριστοδουλόπουλος (Α.Μ. 1093513)}
\institute[CEID, Παν. Πατρών]{
  Τμήμα Μηχανικών Η/Υ \& Πληροφορικής (CEID), Πανεπιστήμιο Πατρών\\
  {\scriptsize \textbf{Επιβλέπων:} Καθηγητής Σπυρίδων Σιούτας}
}
\date{2026}


\begin{document}

% ==============================================================================
% Ενότητα 1: Εισαγωγή & Κίνητρο
% ==============================================================================
% ------------------------------------------------------------------------------
% Slide 1: Τίτλος
% ------------------------------------------------------------------------------
\begin{frame}[plain]
  \vspace{-0.5cm}
  \titlepage
\end{frame}

% ------------------------------------------------------------------------------
% Slide 2: Περιεχόμενα
% ------------------------------------------------------------------------------
\begin{frame}{Δομή Παρουσίασης}
  \tableofcontents
\end{frame}


\section{Εισαγωγή \& Κίνητρο}
% ------------------------------------------------------------------------------
% Slide 3: Κίνητρο & Υπολογιστικό Πρόβλημα
% ------------------------------------------------------------------------------
\begin{frame}{Κίνητρο \& Υπολογιστικό Πρόβλημα}
  \begin{columns}[c]
    \begin{column}{0.56\textwidth}
      \textbf{Αεροδυναμική Σχεδίαση 3D-Printable UAVs:}
      \begin{itemize}
        \item Μη επανδρωμένα αεροσκάφη: μέγιστος λόγος $L/D$.
        \item Αυστηροί περιορισμοί: όγκος ατράκτου, ευστάθεια $C_{m_\alpha} < 0$.
      \end{itemize}

      \vspace{0.1cm}
      \begin{alertblock}{Το Υπολογιστικό Bottleneck}
        \begin{itemize}
          \item Κλασικοί GAs: απαιτούν \textbf{χιλιάδες αξιολογήσεις}.
          \item Κάθε επίλυση VLM/CFD: από εκατοντάδες ms έως λεπτά.
          \item \textbf{Συνολικός χρόνος:} Απαγορευτικός χωρίς HPC clusters.
        \end{itemize}
      \end{alertblock}
    \end{column}
    
    \begin{column}{0.42\textwidth}
      \centering
      \includegraphics[width=\textwidth,height=0.62\textheight,keepaspectratio]{fig1_wing_geometry_3d.png}
      \par\vspace{0.05cm}{\centering\scriptsize 3D Παραμετρική πτέρυγα και άτρακτος.\par}
    \end{column}
  \end{columns}
\end{frame}

% ------------------------------------------------------------------------------
% Slide 4: Στόχος & Επιστημονική Συνεισφορά
% ------------------------------------------------------------------------------
\begin{frame}{Στόχος \& Επιστημονική Συνεισφορά}
  \begin{block}{Κεντρικός Στόχος}
    Δραστική μείωση υπολογιστικού χρόνου \textbf{χωρίς απώλεια ακρίβειας ή σύγκλισης}.
  \end{block}

  \vspace{0.1cm}
  \textbf{Βασικές Επιστημονικές Συνεισφορές:}
  \begin{enumerate}
    \item \textbf{Ιεραρχικός Αξιολογητής 3 Επιπέδων (Multi-Tier $\epsilon$-Bypass):}
    Χωρική μνήμη cache ($\epsilon$-exact), ενεργό surrogate και πλήρης επιλυτής CFD/VLM.
    \item \textbf{Κατανεμημένο Ευρετήριο Μάθησης (LEAD DHT):}
    Chord ring 100 vnodes, προβολή Hilbert 10D $\to$ 1D, και Recursive Model Index (RMI).
    \item \textbf{Αυτόνομη Μικροϋπηρεσία Surrogate σε Rust/C++:}
    GP και MLP με δυναμική επιγραμμική επανεκπαίδευση (online retraining).
    \item \textbf{Κατανεμημένη Αρχιτεκτονική Παραγωγής (Production-Grade Stack):}
    Async Rust (Tokio, Tonic gRPC), Envoy Load Balancing, Kafka Telemetry, 3D Web Monitor.
  \end{enumerate}
\end{frame}


% ==============================================================================
% Ενότητα 2: Αρχιτεκτονική Συστήματος
% ==============================================================================
\section{Αρχιτεκτονική Συστήματος}
% ------------------------------------------------------------------------------
% Slide 5: Συνολική Αρχιτεκτονική
% ------------------------------------------------------------------------------
\begin{frame}{Συνολική Αρχιτεκτονική Συστήματος}
  \begin{columns}[c]
    \begin{column}{0.50\textwidth}
      \textbf{Αυτόνομα Υποσυστήματα gRPC:}
      \begin{itemize}
        \item \textbf{Orchestrators (Rust):} $(\mu+\lambda)$ GA, Gene Store \& pipeline.
        \item \textbf{LEAD DHT Cluster (Rust):} 3 κόμβοι, 60 vnodes, $k$-NN search.
        \item \textbf{Surrogate Node (Rust/HIP):} Fast inference (MLP/GP) \& retraining.
        \item \textbf{Worker Pool (Python):} 3D VLM AeroSandbox.
        \item \textbf{Envoy Proxy:} L7 Load Balancer (Least-Request).
        \item \textbf{Kafka + Fluent-Bit:} Τηλεμετρία \& 3D Real-time Monitor.
      \end{itemize}
    \end{column}

    \begin{column}{0.48\textwidth}
      \centering
      \includegraphics[width=\textwidth,height=0.68\textheight,keepaspectratio]{fig5_system_architecture.png}
      \par\vspace{0.05cm}{\centering\scriptsize Τοπολογία κατανεμημένου συστήματος.\par}
    \end{column}
  \end{columns}
\end{frame}

% ------------------------------------------------------------------------------
% Slide 6: Ιεραρχικός Αξιολογητής Multi-Tier ε-Bypass
% ------------------------------------------------------------------------------
\begin{frame}{Ιεραρχικός Αξιολογητής 3 Επιπέδων (Multi-Tier $\epsilon$-Bypass)}
  \begin{columns}[c]
    \begin{column}{0.54\textwidth}
      \begin{block}{Κριτήρια Διαλογής Ατόμου $\mathbf{x}$}
        \begin{itemize}
          \item \textbf{Tier 1: Χωρική Μνήμη ($\epsilon$-Bypass)}
          \begin{itemize}
            \item $d_{\min} < \epsilon_{\text{exact}} = 10^{-4} \to$ \textbf{Exact Cache Hit}
            \item Άμεση ανάκτηση ($0\ \mu\text{s}$ CPU).
          \end{itemize}
          \vspace{0.05cm}
          \item \textbf{Tier 2: Ενεργό Surrogate (GP/MLP)}
          \begin{itemize}
            \item $d_{\min} \le R$ \textbf{\&} $\sigma(\mathbf{x}) < \theta_{\text{unc}}$
            \item Ταχεία πρόβλεψη ($<1\ \mu\text{s}$).
          \end{itemize}
          \vspace{0.05cm}
          \item \textbf{Tier 3: Πλήρης Επιλυτής VLM}
          \begin{itemize}
            \item Νέα περιοχή ή υψηλή αβεβαιότητα $\sigma \ge \theta$.
            \item gRPC μέσω Envoy σε Python Worker.
            \item Feedback στο LEAD DHT \& Surrogate.
          \end{itemize}
        \end{itemize}
      \end{block}
    \end{column}

    \begin{column}{0.44\textwidth}
      \centering
      \begin{tikzpicture}[node distance=0.65cm, auto,
        block/.style={rectangle, draw=ceidblue, fill=softgray, text width=2.4cm, text centered, rounded corners, minimum height=0.52cm, font=\scriptsize},
        decision/.style={diamond, draw=accentblue, fill=accentblue!10, text width=1.8cm, text centered, inner sep=0pt, font=\tiny},
        line/.style={draw=ceidblue, -stealth', thick}]
        
        \node [block, fill=ceidblue!10] (input) {Υποψήφιο $\mathbf{x}$};
        \node [decision, below=0.35cm of input] (tier1) {$d_{\min} \le \epsilon$;};
        \node [block, right=0.3cm of tier1, fill=green!20] (t1_res) {Tier 1: Cache ($0$ cost)};
        \node [decision, below=0.45cm of tier1] (tier2) {$d \le R \land \sigma < \theta$;};
        \node [block, right=0.3cm of tier2, fill=yellow!25] (t2_res) {Tier 2: Surrogate};
        \node [block, below=0.45cm of tier2, fill=red!20] (t3_res) {Tier 3: Worker VLM};
        
        \path [line] (input) -- (tier1);
        \path [line] (tier1) -- node [above] {\tiny Ναι} (t1_res);
        \path [line] (tier1) -- node [right] {\tiny Όχι} (tier2);
        \path [line] (tier2) -- node [above] {\tiny Ναι} (t2_res);
        \path [line] (tier2) -- node [right] {\tiny Όχι} (t3_res);
      \end{tikzpicture}
    \end{column}
  \end{columns}
\end{frame}

% ------------------------------------------------------------------------------
% Slide 7: Μοντέλο Νήσων & Ασύγχρονη Μετανάστευση
% ------------------------------------------------------------------------------
\begin{frame}{Μοντέλο Νήσων \& Ασύγχρονη Μετανάστευση}
  \begin{columns}[c]
    \begin{column}{0.50\textwidth}
      \textbf{Αποκεντρωμένη Εξέλιξη:}
      \begin{itemize}
        \item Orchestrators σε \textbf{τοπολογία δακτυλίου (Ring)}.
        \item $(\mu + \lambda)$ GA ανά νήσο ($P_{\text{local}} = 100$).
        \item \textbf{Αποφυγή πρόωρης σύγκλισης:} Παράλληλη εξερεύνηση διαφορετικών περιοχών 10D χώρου.
      \end{itemize}

      \vspace{0.1cm}
      \textbf{Πρωτόκολλο Μετανάστευσης:}
      \begin{itemize}
        \item Ανά διάστημα $\tau = 5$ γενεών: ανταλλαγή των $M_c = 3$ ελίτ ατόμων στον επόμενο γείτονα.
        \item Ασύγχρονη gRPC κλήση (\texttt{MigrateElites}).
      \end{itemize}
    \end{column}

    \begin{column}{0.48\textwidth}
      \centering
      \includegraphics[width=\textwidth,height=0.62\textheight,keepaspectratio]{fig1_single_vs_multi_island.png}
      \par\vspace{0.05cm}{\centering\scriptsize Σύγκριση Μονής Νήσου έναντι Δακτυλίου Νήσων.\par}
    \end{column}
  \end{columns}
\end{frame}


% ==============================================================================
% Ενότητα 3: LEAD DHT & Χωρική Ευρετηρίαση
% ==============================================================================
\section{LEAD DHT \& Χωρική Ευρετηρίαση}
% ------------------------------------------------------------------------------
% Slide 8: Hilbert Space-Filling Curves (10D -> 1D)
% ------------------------------------------------------------------------------
\begin{frame}{Χωρική Προβολή Hilbert (10D $\to$ 1D Locality)}
  \begin{columns}[c]
    \begin{column}{0.52\textwidth}
      \textbf{Η Πρόκληση των 10 Διαστάσεων:}
      \begin{itemize}
        \item Συμβατικά DHTs (SHA-1): καταστρέφουν τη γειτνίαση.
        \item Αδύνατη η εκτέλεση $k$-NN σε κλασικά P2P.
      \end{itemize}

      \vspace{0.1cm}
      \textbf{Καμπύλη Hilbert ($p=6$):}
      \begin{itemize}
        \item Χαρτογράφηση $\mathbb{R}^{10} \to \text{uint64}$ με διατήρηση εγγύτητας.
        \item \textbf{Multi-Probe ($C=3$ περιστροφές):} Εξάλειψη ασυνεχειών στα όρια κυψελών.
        \item Διπλασιασμός συσχέτισης ($r = 0.333$) και $6.5\times$ μείωση της στρέβλωσης $P_{90}$.
      \end{itemize}
    \end{column}

    \begin{column}{0.46\textwidth}
      \centering
      \includegraphics[width=\textwidth,height=0.65\textheight,keepaspectratio]{fig5_hilbert_spatial_heatmap.png}
      \par\vspace{0.05cm}{\centering\scriptsize 10D vs 1D Hilbert Distance Heatmap.\par}
    \end{column}
  \end{columns}
\end{frame}

% ------------------------------------------------------------------------------
% Slide 9: Δακτύλιος Chord & Learned Index (RMI)
% ------------------------------------------------------------------------------
\begin{frame}{Δακτύλιος Chord \& Learned Index (RMI)}
  \begin{columns}[c]
    \begin{column}{0.50\textwidth}
      \textbf{Κατανεμημένος Δακτύλιος LEAD:}
      \begin{itemize}
        \item 100 Virtual Nodes (vnodes) ανά κόμβο.
        \item Δρομολόγηση $O(\log N)$ hops με finger tables.
      \end{itemize}

      \vspace{0.1cm}
      \textbf{Ευρετήριο RMI (Recursive Model Index):}
      \begin{itemize}
        \item Αντικατάσταση B-Tree με μοντέλο 2 επιπέδων (Root Linear $\to$ 256 Leaf Linear Models).
        \item \textbf{Αποτέλεσμα:} Χρόνος ανάκτησης $\approx 41\text{ ns}$ ($1.81\times$ ταχύτερο του B-Tree).
        \item \textbf{$3\times$ χαμηλότερη μνήμη RAM} ($785\text{ KB}$ vs $2.34\text{ MB}$).
      \end{itemize}
    \end{column}

    \begin{column}{0.48\textwidth}
      \centering
      \includegraphics[width=\textwidth,height=0.65\textheight,keepaspectratio]{fig3_rmi_vs_btree_benchmark.png}
      \par\vspace{0.05cm}{\centering\scriptsize RMI vs std::BTreeMap Latency.\par}
    \end{column}
  \end{columns}
\end{frame}


% ==============================================================================
% Ενότητα 4: Μοντέλα Surrogate & Αεροδυναμική VLM
% ==============================================================================
\section{Μοντέλα Surrogate \& Αεροδυναμική VLM}
% ------------------------------------------------------------------------------
% Slide 10: Μοντέλα Υποκατάστασης (Surrogate Node)
% ------------------------------------------------------------------------------
\begin{frame}{Μοντέλα Υποκατάστασης (Surrogate Node)}
  \begin{columns}[c]
    \begin{column}{0.54\textwidth}
      \textbf{Απαιτήσεις \& Σχεδιασμός:}
      \begin{itemize}
        \item Χαμηλή καθυστέρηση πρόβλεψης ($< 1\text{ ms}$).
        \item Εκτίμηση αβεβαιότητας $\sigma(\mathbf{x})$ για ασφαλές gate.
        \item Online επανεκπαίδευση (sliding window 500).
      \end{itemize}

      \vspace{0.08cm}
      \textbf{Συγκριτική Αξιολόγηση:}
      \begin{itemize}
        \item \textbf{Gaussian Process (Matérn 5/2 ARD):}\\
        Υψηλή ακρίβεια ($R^2 = 0.879$) \& αβεβαιότητα $\sigma(\mathbf{x})$.
        \item \textbf{MLP Neural Network ($64 \times 32$):}\\
        Εξαιρετική ταχύτητα ($>1.2 \times 10^6\text{ evals/s}$).
      \end{itemize}
    \end{column}

    \begin{column}{0.44\textwidth}
      \centering
      \includegraphics[width=\textwidth,height=0.60\textheight,keepaspectratio]{fig2_surrogate_pred_vs_true.png}
      \par\vspace{0.05cm}{\centering\scriptsize Parity Plot: Πραγματικό vs Προβλεπόμενο $L/D$.\par}
    \end{column}
  \end{columns}
\end{frame}

% ------------------------------------------------------------------------------
% Slide 11: Ανάλυση Ευαισθησίας Παραμέτρων (ARD GP)
% ------------------------------------------------------------------------------
\begin{frame}{Ανάλυση Αεροδυναμικής Ευαισθησίας (ARD)}
  \begin{columns}[c]
    \begin{column}{0.52\textwidth}
      \textbf{Automatic Relevance Determination:}
      \begin{itemize}
        \item Μήκος κλίμακας $\ell_d$: μικρό $\ell_d \to$ υψηλή ευαισθησία.
      \end{itemize}

      \vspace{0.1cm}
      \textbf{Κυρίαρχες Παράμετροι Σχεδίασης:}
      \begin{enumerate}
        \item \texttt{naca\_m} (Καμπυλότητα): \textbf{40.3\%} ($\ell = 0.83$).
        \item \texttt{wing\_twist} (Συστροφή): \textbf{16.2\%} ($\ell = 2.06$).
        \item \texttt{naca\_t} (Πάχος): \textbf{16.2\%} ($\ell = 2.07$).
        \item \texttt{fuse\_diam} (Διάμετρος): \textbf{12.9\%} ($\ell = 2.59$).
      \end{enumerate}
      \vspace{0.05cm}
      {\footnotesize $\to$ Πλήρης συμφωνία με την αεροδυναμική θεωρία.}
    \end{column}

    \begin{column}{0.46\textwidth}
      \centering
      \includegraphics[width=\textwidth,height=0.65\textheight,keepaspectratio]{fig2_surrogate_parameter_sensitivity.png}
      \par\vspace{0.05cm}{\centering\scriptsize Σχετική βαρύτητα και Length Scales παραμέτρων.\par}
    \end{column}
  \end{columns}
\end{frame}

% ------------------------------------------------------------------------------
% Slide 12: Φυσικός Εξομοιωτής VLM Worker (Tier 3)
% ------------------------------------------------------------------------------
\begin{frame}{Φυσικός Εξομοιωτής Worker (Tier 3)}
  \begin{columns}[c]
    \begin{column}{0.52\textwidth}
      \textbf{AeroSandbox 3D VLM:}
      \begin{itemize}
        \item Επίλυση δυναμικού πεδίου ροής πτέρυγας \& ατράκτου.
        \item Άντωση $C_L$, οπισθέλκουσα $C_D$, ευστάθεια $C_{m_\alpha} < 0$.
      \end{itemize}

      \vspace{0.1cm}
      \begin{block}{Συνάρτηση Καταλληλότητας (Fitness)}
        \vspace{-0.1cm}
        {\footnotesize
        \[
        F(\mathbf{x}) = \begin{cases}
          \frac{L}{D} - \sum \mathcal{P}, & \text{εφικτή} \\
          -10^9, & \text{μη εφικτή}
        \end{cases}
        \]
        }
        \vspace{-0.25cm}
        {\scriptsize Ποινές $\mathcal{P}$: αεροδυναμική, όγκος ($V \ge 2.5 \cdot 10^5\text{ mm}^3$), ευστάθεια.}
      \end{block}
    \end{column}

    \begin{column}{0.46\textwidth}
      \centering
      \includegraphics[width=\textwidth,height=0.65\textheight,keepaspectratio]{fig5_aero_pareto_radar.png}
      \par\vspace{0.05cm}{\centering\scriptsize Radar Chart: Σύγκριση αεροδυναμικών δεικτών.\par}
    \end{column}
  \end{columns}
\end{frame}


% ==============================================================================
% Ενότητα 5: Πειραματική Αξιολόγηση
% ==============================================================================
\section{Πειραματική Αξιολόγηση}
% ------------------------------------------------------------------------------
% Slide 13: Σύγκλιση Καταλληλότητας & Ευστάθεια
% ------------------------------------------------------------------------------
\begin{frame}{Σύγκλιση Καταλληλότητας \& Αναπαραγωγιμότητα}
  \begin{columns}[c]
    \begin{column}{0.50\textwidth}
      \textbf{Δυναμική Σύγκλισης (47 Γενεές):}
      \begin{itemize}
        \item Ταχεία άνοδος (από $L/D = 9.37 \to 11.68$).
        \item Ασυμπτωτική σύγκλιση στο \textbf{$L/D \approx 12.28$}.
        \item \textbf{Βελτίωση $+17.8\%$} έναντι του σχεδίου βάσης.
      \end{itemize}

      \vspace{0.1cm}
      \textbf{Αξιοπιστία Πολλαπλών Seeds:}
      \begin{itemize}
        \item 5 ανεξάρτητα πειράματα (Seeds 42, 101, 777, 999, 2024).
        \item Όλα συγκλίνουν στο $12.27 \pm 0.02$.
        \item Μηδενική ευαισθησία σε αρχικοποίηση.
      \end{itemize}
    \end{column}

    \begin{column}{0.48\textwidth}
      \centering
      \includegraphics[width=\textwidth,height=0.65\textheight,keepaspectratio]{fig5_fitness_multi_seed_ribbon.png}
      \par\vspace{0.05cm}{\centering\scriptsize Κορδέλα σύγκλισης $\pm 1\sigma$ για 5 seeds.\par}
    \end{column}
  \end{columns}
\end{frame}

% ------------------------------------------------------------------------------
% Slide 14: Ανάλυση Bypass Ratio & Επιτάχυνση
% ------------------------------------------------------------------------------
\begin{frame}{Ανάλυση Bypass Ratio \& Επιτάχυνση Χρόνου}
  \begin{columns}[c]
    \begin{column}{0.50\textwidth}
      \begin{block}{Ρυθμός Παράκαμψης (Bypass Ratio)}
        $$BR = \frac{N_{\text{Tier 1}} + N_{\text{Tier 2}}}{N_{\text{total}}} = 1 - \frac{N_{\text{CFD}}}{N_{\text{total}}}$$
      \end{block}
      \begin{itemize}
        \item \textbf{Warmup Phase ($t \le 5$):} $BR < 15\%$ (συλλογή δειγμάτων).
        \item \textbf{Μόνιμη Κατάσταση ($t \ge 15$):} \textbf{$BR > 75\%$}.
        \item \textbf{Μείωση κλήσεων VLM:} από 4.500 σε 1.972 ($-56.2\%$).
        \item \textbf{Wall-Clock Speedup: $2.28\times$} (6.0 min vs 13.8 min).
      \end{itemize}
    \end{column}

    \begin{column}{0.48\textwidth}
      \centering
      \includegraphics[width=\textwidth,height=0.65\textheight,keepaspectratio]{fig2_tier_breakdown_percent.png}
      \par\vspace{0.05cm}{\centering\scriptsize Ποσοστό συμμετοχής Tiers ανά γενιά.\par}
    \end{column}
  \end{columns}
\end{frame}

% ------------------------------------------------------------------------------
% Slide 15: Wall-Clock Speedup Breakdown
% ------------------------------------------------------------------------------
\begin{frame}{Συγκριτική Επιτάχυνση Pipeline (Wall-Clock)}
  \begin{columns}[c]
    \begin{column}{0.48\textwidth}
      \centering
      \includegraphics[width=\textwidth,height=0.65\textheight,keepaspectratio]{fig2_pipeline_speedup_barchart.png}
      \par\vspace{0.05cm}{\centering\scriptsize Χρόνος εκτέλεσης ανά διαμόρφωση pipeline.\par}
    \end{column}

    \begin{column}{0.50\textwidth}
      \begin{table}
        \centering
        \scriptsize
        \begin{tabular}{lccc}
          \toprule
          \textbf{Διαμόρφωση} & \textbf{Χρόνος} & \textbf{Speedup} & \textbf{Best $F$} \\
          \midrule
          1-Tier (CFD Only) & 13.8 min & $1.00\times$ & 12.27 \\
          2-Tier (Cache+CFD) & 10.9 min & $1.27\times$ & 12.27 \\
          \textbf{3-Tier (Full)} & \textbf{6.0 min} & \textbf{2.28$\times$} & \textbf{12.27} \\
          \bottomrule
        \end{tabular}
      \end{table}

      \vspace{0.1cm}
      \begin{exampleblock}{Μηδενικός Συμβιβασμός Ποιότητας}
        \footnotesize
        Η επιτάχυνση $2.28\times$ επιτυγχάνεται \textbf{χωρίς καμία απώλεια ακρίβειας} ($F = 12.27$ σε όλες τις περιπτώσεις).
      \end{exampleblock}
    \end{column}
  \end{columns}
\end{frame}

% ------------------------------------------------------------------------------
% Slide 16: Κλιμακωσιμότητα Workers & Εξισορρόπηση Envoy
% ------------------------------------------------------------------------------
\begin{frame}{Κλιμακωσιμότητα Workers \& Εξισορρόπηση Envoy}
  \begin{columns}[c]
    \begin{column}{0.50\textwidth}
      \textbf{Οριζόντια Κλιμάκωση:}
      \begin{itemize}
        \item Παράλληλο κλάσμα $p \approx 0.96$ (Amdahl / Gustafson).
        \item Στους 8 workers: επιτάχυνση $7.15\times$ ($\eta = 89.4\%$).
        \item Στους 32 workers: \textbf{$21.54\times$} ($98.2\ \text{άτομα/s}$).
      \end{itemize}

      \vspace{0.1cm}
      \textbf{Envoy Least-Request vs Round-Robin:}
      \begin{itemize}
        \item Εξάλειψη head-of-line blocking.
        \item Μείωση του $P_{99}$ latency κατά $28\%$.
      \end{itemize}
    \end{column}

    \begin{column}{0.48\textwidth}
      \centering
      \includegraphics[width=\textwidth,height=0.65\textheight,keepaspectratio]{fig4_worker_scalability.png}
      \par\vspace{0.05cm}{\centering\scriptsize Κλιμακωσιμότητα Workers ($W=1\dots 32$).\par}
    \end{column}
  \end{columns}
\end{frame}

% ------------------------------------------------------------------------------
% Slide 17: Επιτάχυνση Υλικού AMD ROCm GPU & Benchmarks
% ------------------------------------------------------------------------------
\begin{frame}{Επιτάχυνση Υλικού AMD ROCm GPU (10 Trials)}
  \vspace{-0.2cm}
  \begin{columns}[T]
    \begin{column}{0.32\textwidth}
      \begin{block}{GP Covariance ($N=10^3$)}
        \scriptsize
        \textbf{ROCm:} $\mathbf{0.64 \pm 0.01\text{ ms}}$\\
        \textbf{CPU (12t):} $1.23 \pm 0.12\text{ ms}$\\
        $\mathbf{1.91\times}$ speedup ($\sigma \approx 0.01\text{ ms}$)
      \end{block}
    \end{column}
    \begin{column}{0.34\textwidth}
      \begin{block}{Cholesky ($LL^T = K$)}
        \scriptsize
        \textbf{CPU ($N=10^3$):} $\mathbf{103.9 \pm 0.5\text{ ms}}$\\
        $\mathcal{O}(N^3)$ κυβικό bottleneck\\
        $\to$ Αιτιολόγηση MLP στο Tier 2
      \end{block}
    \end{column}
    \begin{column}{0.32\textwidth}
      \begin{block}{MLP Scaling ($M=10^4$)}
        \scriptsize
        \textbf{ROCm:} $\mathbf{0.46 \pm 0.01\text{ ms}}$ ($\mathbf{4.12\times}$)\\
        \textbf{Throughput:} $\mathbf{21.6 \times 10^6\text{ ev/s}}$\\
        \textit{Hilbert: $4.4\ \mu\text{s}$ | Store: $18\text{ ns}$}
      \end{block}
    \end{column}
  \end{columns}

  \vspace{0.1cm}
  \begin{center}
    \includegraphics[width=0.98\textwidth,height=0.46\textheight,keepaspectratio]{fig2_surrogate_cuda_vs_cpu.png}
    \par\vspace{0.03cm}{\tiny\textcolor{darktext!70}{Host CPU vs AMD ROCm GPU (RX 7600 XT). Σκιώδεις ζώνες: $\pm 1\sigma$ εμπειρική αβεβαιότητα (10 ανεξάρτητες δοκιμές).}\par}
  \end{center}
\end{frame}


% ==============================================================================
% Ενότητα 6: Συμπεράσματα & Μελλοντική Έρευνα
% ==============================================================================
\section{Συμπεράσματα}
% ------------------------------------------------------------------------------
% Slide 18: Συμπεράσματα
% ------------------------------------------------------------------------------
\begin{frame}{Συμπεράσματα}
  \begin{block}{Κύρια Επιτεύγματα της Εργασίας}
    \begin{enumerate}
      \item \textbf{Δραστική Μείωση Κόστους:} $BR = 75.5\%$ και συνολική επιτάχυνση \textbf{$2.28\times$} σε πραγματικό χρόνο.
      \item \textbf{Αεροδυναμική Βελτιστοποίηση:} Αύξηση λόγου $L/D$ από $11.07 \to \mathbf{13.04}$ (\textbf{$+17.8\%$}) με πλήρη ικανοποίηση περιορισμών.
      \item \textbf{Πρωτότυπη Χωρική Ευρετηρίαση:} Καμπύλες Hilbert multi-probe με Learned Index (RMI) σε Chord DHT: lookups $41\text{ ns}$, $3\times$ λιγότερη RAM.
      \item \textbf{Αρχιτεκτονική Παραγωγής σε Rust:} Μηδενικά memory leaks, κλιμάκωση έως 32 κατανεμημένους εργάτες ($21.5\times$ speedup).
    \end{enumerate}
  \end{block}
\end{frame}

% ------------------------------------------------------------------------------
% Slide 19: Μελλοντική Έρευνα
% ------------------------------------------------------------------------------
\begin{frame}{Μελλοντικές Επεκτάσεις}
  \begin{itemize}
    \item \textbf{Επιλυτές Υψηλότερης Πιστότητας (CFD RANS):}
    Υβριδική σύζευξη VLM με OpenFOAM / SU2 σε GPU συστοιχίες.
    
    \vspace{0.15cm}
    \item \textbf{Federated Learning στα Μοντέλα Surrogate:}
    Ασύγχρονος διαμοιρασμός βαρών (FedAvg) των MLP surrogates μέσω Kafka.
    
    \vspace{0.15cm}
    \item \textbf{Πολυ-Κριτηριακή Εξέλιξη NSGA-III:}
    Απευθείας εύρεση τρισδιάστατου μετώπου Pareto ($L/D$, Βάρος, Payload, Κόστος).
    
    \vspace{0.15cm}
    \item \textbf{Κατασκευή \& Πτητική Δοκιμή:}
    3D εκτύπωση της βέλτιστης γεωμετρίας και πειραματική επικύρωση σε αεροσήραγγα.
  \end{itemize}
\end{frame}

% ------------------------------------------------------------------------------
% Slide 20: Τέλος Παρουσίασης & Q&A
% ------------------------------------------------------------------------------
\begin{frame}[standout]
  \Huge Ευχαριστώ!

  \vspace{0.6cm}
  \normalsize Ερωτήσεις \& Συζήτηση

  \vspace{0.8cm}
  \footnotesize
  Ευστάθιος Παναγιώτης Χριστοδουλόπουλος\\
  Τμήμα Μηχανικών Η/Υ \& Πληροφορικής, Πανεπιστήμιο Πατρών\\
  \texttt{up1093513@ac.upatras.gr}
\end{frame}


\end{document}
